\documentclass[11pt]{article}
\usepackage[utf8]{inputenc}
\usepackage{amsmath,amssymb}
\usepackage[margin=1in]{geometry}
\usepackage{hyperref}
\emergencystretch=3em

\title{The product-kernel block with unequal exponents}
\author{Claude, with Daniel Murfet}
\date{2026-09-07}

\begin{document}
\maketitle
\sloppy

\begin{abstract}
For $p_1=(h_1+1)/k_1<p_2=(h_2+1)/k_2$ and a nonnegative kernel $g$, the product-kernel block
$J(t)=\int\!\!\int_{(0,b]^2}u^{h_1}v^{h_2}g(tu^{k_1}v^{k_2})$ is bounded by the smaller exponent
alone (Astra~\#5(d)):
\[
J(t)\le\frac{b^{k_2(p_2-p_1)}}{k_1k_2(p_2-p_1)}\,M_{p_1}(g)\,t^{-p_1},\qquad t>0 .
\]
The density of unit~103 with constant amplitude is
$k_2^{-1}r^{p_2-1}g(tr)(\varepsilon^{-\gamma}-b^{-\gamma})/\gamma$ with $\gamma=k_1(p_2-p_1)$ and
$\varepsilon=(r/b^{k_2})^{1/k_1}$, which is at most
$b^{k_2(p_2-p_1)}/(k_1k_2(p_2-p_1))\cdot r^{p_1-1}g(tr)$. Together with the equal-exponent bound of
unit~111 this makes the mixed remainder of the $d=2$ cutoff theorem fully quantitative in both regimes.
Zero \texttt{sorry}/\texttt{axiom}.
\end{abstract}

\section{The things formalised}
{\small
\begin{verbatim}
theorem genDensity_const_le (b p₁ p₂ : ℝ) (h₁ k₁ k₂ : ℕ) (hb : 0 < b) (hk₁ : 0 < k₁)
    (hk₂ : 0 < k₂) (hp₁ : ((h₁ : ℝ) + 1) / k₁ = p₁) (hlt : p₁ < p₂) (g : ℝ → ℝ)
    (hg0 : ∀ s, 0 ≤ g s) (r s : ℝ) (hr : 0 < r) (hrR : r ≤ b ^ (k₁ + k₂)) :
    genDensity p₂ b h₁ k₁ k₂ (fun _ s => g s) r s
      ≤ b ^ ((k₂ : ℝ) * (p₂ - p₁)) / ((k₁ : ℝ) * k₂ * (p₂ - p₁)) * (r ^ (p₁ - 1) * g s)
theorem prodKernel_le_unequal (b p₁ p₂ : ℝ) (h₁ h₂ k₁ k₂ : ℕ) (hb : 0 < b) (hk₁ : 0 < k₁)
    (hk₂ : 0 < k₂) (hp₁ : ((h₁ : ℝ) + 1) / k₁ = p₁) (hp₂ : ((h₂ : ℝ) + 1) / k₂ = p₂)
    (hlt : p₁ < p₂) (g : ℝ → ℝ) (hg : Continuous g) (hg0 : ∀ s, 0 ≤ g s)
    (hM : IntegrableOn (fun s => s ^ (p₁ - 1) * g s) (Ioi 0)) (t : ℝ) (ht : 0 < t) :
    prodKernel b h₁ h₂ k₁ k₂ g t
      ≤ b ^ ((k₂ : ℝ) * (p₂ - p₁)) / ((k₁ : ℝ) * k₂ * (p₂ - p₁)) * kernelMoment p₁ g
        * t ^ (-p₁)
\end{verbatim}
}

\section{Mathematics}

With weight exponent $w=h_1-k_1p_2$, $w+1=-k_1(p_2-p_1)=-\gamma<0$, so
$\int_\varepsilon^bu^w\,du=(\varepsilon^{-\gamma}-b^{-\gamma})/\gamma\le\varepsilon^{-\gamma}/\gamma$
and $\varepsilon^{-\gamma}=(r/b^{k_2})^{-(p_2-p_1)}=r^{-(p_2-p_1)}b^{k_2(p_2-p_1)}$; the powers of $r$
combine to $r^{p_2-1-(p_2-p_1)}=r^{p_1-1}$. The block integral is the $r$-integral of the density over
$(0,R]$ (unit~103 at $\beta=0$); \texttt{integral\_mono\_of\_nonneg} compares it with the dominating
integrand without needing integrability of the density itself (only nonnegativity and integrability
of the bound, which follows from boundedness of $g$ on $[0,tR]$), the substitution $x=tr$ produces
$t^{-p_1}$, and the resulting moment over $(0,tR]$ is at most $M_{p_1}$.

\section{Paper-facing meaning}

The mixed remainder of the rectangular decomposition has shifted exponents
$(h_1+M_1+1)/k_1$ and $(h_2+M_2+1)/k_2$, which coincide only when $M_1/k_1=M_2/k_2$; the present
bound covers the generic unequal case with an explicit constant, the unit-111 bound the equal case.
The $1/(p_2-p_1)$ factor is the artificial blow-up Astra~\#5(d) warned about near a collision; the
gap-uniform alternative $(1-x^\delta)/\delta\le\log(1/x)$ remains a follow-up.

\section{Lean notes}

\begin{itemize}
\item \texttt{integral\_mono\_of\_nonneg} needs integrability of the \emph{dominating} function
only; use it whenever the smaller integrand is awkward to show integrable.
\item Power identities: after \texttt{Real.div\_rpow}, \texttt{← Real.rpow\_natCast},
\texttt{← Real.rpow\_mul}, \texttt{div\_eq\_mul\_inv} and \texttt{← Real.rpow\_neg}, finish with
\texttt{congr 2; ring} on the exponent; \texttt{rw} cannot see through
\texttt{(b \^{} x)⁻¹ = b \^{} (-x)} in the other direction once other rewrites have fired.
\item \texttt{div\_neg}, \texttt{← neg\_div}, \texttt{neg\_sub} in that order rewrite
$(a-c)/(-d)$ to $(c-a)/d$.
\item \texttt{field\_simp; try ring} accommodates both outcomes (closed or not).
\end{itemize}

\end{document}
